\documentclass[12pt,a4paper]{report}

% ==========================================
% PACKAGES
% ==========================================
\usepackage[utf8]{inputenc}
\usepackage[T1]{fontenc}
\usepackage[francais]{babel}
\usepackage{geometry}
\usepackage{xcolor}
\usepackage{booktabs}
\usepackage{tabularx}
\usepackage{longtable}
\usepackage{array}
\usepackage{graphicx}
\usepackage{listings}
\usepackage{fancyhdr}
\usepackage{titlesec}
\usepackage{tcolorbox}
\usepackage{hyperref}
\usepackage{enumitem}
\usepackage{float}

\geometry{
  headheight=15pt,
  top=2.5cm, bottom=2.5cm,
  left=2.5cm, right=2.5cm
}

% ==========================================
% COULEURS GLOBALES
% ==========================================
\definecolor{primary}{RGB}{30,56,100}
\definecolor{secondary}{RGB}{46,117,182}
\definecolor{codegray}{RGB}{150,150,150}
\definecolor{accent}{RGB}{255,165,0}
\definecolor{lightblue}{RGB}{235,245,255}
\definecolor{successgreen}{RGB}{0,128,0}

% ==========================================
% COULEURS ET STYLE DES LISTINGS (CODE)
% ==========================================
\definecolor{lightgray}{rgb}{0.96,0.96,0.96}
\definecolor{darkgray}{rgb}{0.3,0.3,0.3}
\definecolor{codegreen}{RGB}{0,128,0}
\definecolor{codepurple}{RGB}{128,0,128}
\definecolor{codeblue}{RGB}{0,0,180}

\lstdefinestyle{shell}{
  backgroundcolor=\color{lightgray},
  frame=single,
  rulecolor=\color{darkgray},
  basicstyle=\ttfamily\small,
  breaklines=true,
  numbers=left,
  numberstyle=\tiny\color{darkgray},
  numbersep=8pt,
  showstringspaces=false,
  tabsize=2,
  aboveskip=1em,
  belowskip=1em
}

\lstdefinestyle{javastyle}{
  backgroundcolor=\color{lightgray},
  commentstyle=\color{codegreen}\itshape,
  keywordstyle=\color{codeblue}\bfseries,
  stringstyle=\color{codepurple},
  basicstyle=\ttfamily\small,
  breaklines=true,
  frame=single,
  rulecolor=\color{darkgray},
  numbers=left,
  numberstyle=\tiny\color{darkgray},
  numbersep=8pt,
  tabsize=4,
  showstringspaces=false,
  aboveskip=1em,
  belowskip=1em,
  captionpos=b
}

\lstdefinestyle{sqlstyle}{
  backgroundcolor=\color{lightgray},
  keywordstyle=\color{codeblue}\bfseries,
  stringstyle=\color{codepurple},
  basicstyle=\ttfamily\small,
  breaklines=true,
  frame=single,
  rulecolor=\color{darkgray},
  numbers=left,
  numberstyle=\tiny\color{darkgray},
  numbersep=8pt,
  tabsize=4,
  showstringspaces=false,
  language=SQL,
  aboveskip=1em,
  belowskip=1em,
  captionpos=b
}

\lstdefinestyle{yamlstyle}{
  backgroundcolor=\color{lightgray},
  keywordstyle=\color{codeblue}\bfseries,
  stringstyle=\color{codepurple},
  basicstyle=\ttfamily\small,
  breaklines=true,
  frame=single,
  rulecolor=\color{darkgray},
  numbers=left,
  numberstyle=\tiny\color{darkgray},
  numbersep=8pt,
  tabsize=2,
  showstringspaces=false,
  aboveskip=1em,
  belowskip=1em,
  captionpos=b
}

\lstdefinestyle{graphqlstyle}{
  backgroundcolor=\color{lightgray},
  keywordstyle=\color{codeblue}\bfseries,
  stringstyle=\color{codepurple},
  basicstyle=\ttfamily\small,
  breaklines=true,
  frame=single,
  rulecolor=\color{darkgray},
  numbers=left,
  numberstyle=\tiny\color{darkgray},
  numbersep=8pt,
  tabsize=4,
  showstringspaces=false,
  aboveskip=1em,
  belowskip=1em,
  captionpos=b
}

% ==========================================
% STYLE DES TITRES
% ==========================================
\titleformat{\chapter}[display]
  {\normalfont\huge\bfseries\color{primary}}
  {\chaptertitlename\ \thechapter}{20pt}{\Huge}
\titleformat{\section}
  {\normalfont\Large\bfseries\color{secondary}}
  {\thesection}{1em}{}
\titleformat{\subsection}
  {\normalfont\large\bfseries\color{primary}}
  {\thesubsection}{1em}{}
\titleformat{\subsubsection}
  {\normalfont\normalsize\bfseries}
  {\thesubsubsection}{1em}{}

% ==========================================
% EN-TÊTE ET PIED DE PAGE
% ==========================================
\pagestyle{fancy}
\fancyhf{}
\fancyhead[L]{\small\color{secondary}Gestion de Flotte de V\'ehiculess}
\fancyhead[R]{\small\color{secondary}Master 1 GIL --- 2025/2026}
\fancyfoot[C]{\thepage}
\renewcommand{\headrulewidth}{0.4pt}

% ==========================================
% BOÎTES ADR ET INFO
% ==========================================
\tcbuselibrary{skins,breakable}
\newtcolorbox{adrbox}[2][]{
  enhanced,
  colback=lightblue,
  colframe=secondary,
  fonttitle=\bfseries\color{white},
  title=#2,
  #1
}

\newtcolorbox{infobox}{
  enhanced,
  colback=yellow!10,
  colframe=accent,
  fonttitle=\bfseries,
}

% ==========================================
\begin{document}

% ==========================================
% PAGE DE TITRE
% ==========================================
\begin{titlepage}
\centering
\vspace*{2cm}

{\color{primary}\rule{\linewidth}{1pt}}
\vspace{0.5cm}

{\Huge\bfseries\color{primary} Rapport Technique\par}
\vspace{0.4cm}
{\LARGE\color{secondary} Service de Gestion des V\'ehiculess\par}
\vspace{0.3cm}
{\large Microservice Spring Boot --- CRUD, Kafka, GraphQL, Tests et Observabilit\'e\par}

\vspace{0.5cm}
{\color{primary}\rule{\linewidth}{1pt}}

\vspace{2cm}

{\Large Syst\`eme de Gestion de Flotte de V\'ehiculess\par}
\vspace{1cm}
{\large Master 1 GIL --- 2025/2026\par}
\vspace{0.5cm}
{\large Mars 2026\par}

\vfill
\end{titlepage}

% ============================================================
\chapter{Introduction}
% ============================================================

Ce rapport d\'ecrit en d\'etail l'impl\'ementation du microservice \textbf{vehicule-service}, un composant central du syst\`eme de gestion de flotte de v\'ehicules. Ce service g\`ere le cycle de vie complet des v\'ehicules : cr\'eation, consultation, mise \`a jour et suppression (CRUD), tout en int\'egrant des technologies modernes pour la communication asynchrone (Kafka), les API flexibles (GraphQL), l'observabilit\'e (OpenTelemetry) et les tests automatis\'es.

\section{Architecture g\'en\'erale}

Le service suit une architecture en couches classique :

\begin{enumerate}
    \item \textbf{Couche Contr\^oleur} --- Expose les API REST et GraphQL
    \item \textbf{Couche Service} --- Contient la logique m\'etier
    \item \textbf{Couche Repository} --- G\`ere la persistance avec Spring Data JPA
    \item \textbf{Couche \'Ev\'enementielle} --- Publie et consomme des \'ev\'enements via Kafka
\end{enumerate}

\begin{infobox}
\textbf{Port du service :} 8081 \\
\textbf{Base de donn\'ees :} PostgreSQL (fleet\_vehicles) \\
\textbf{Messagerie :} Apache Kafka (localhost:9092) \\
\textbf{Tracing :} OpenTelemetry (OTLP sur localhost:4317)
\end{infobox}

% ============================================================
\chapter{Interface REST et Alignement OpenAPI}
% ============================================================

Le service expose une interface REST align\'ee sur la sp\'ecification OpenAPI de la passerelle (\texttt{gateway/openapi/vehicule-openapi.yaml}). Les points d'entr\'ee utilisent le nommage en fran\c{c}ais (\texttt{/api/vehicules}) et le champ identifiant est \texttt{id\_vehicule}.

\section{Points d'entr\'ee (Endpoints)}

\begin{adrbox}{Endpoints REST --- Base path: /api/vehicules}
\begin{tabularx}{\textwidth}{l l X}
\toprule
\textbf{M\'ethode} & \textbf{Endpoint} & \textbf{Description} \\
\midrule
\texttt{POST} & \texttt{/} & Cr\'eation d'un v\'ehicule \\
\texttt{GET} & \texttt{/} & Liste de tous les v\'ehicules \\
\texttt{GET} & \texttt{/disponibles} & Liste uniquement les v\'ehicules disponibles \\
\texttt{GET} & \texttt{/\{id\}} & D\'etail d'un v\'ehicule par UUID \\
\texttt{PUT} & \texttt{/\{id\}} & Mise \`a jour compl\`ete \\
\texttt{PATCH} & \texttt{/\{id\}/statut} & Mise \`a jour du statut (Kafka) \\
\texttt{PATCH} & \texttt{/\{id\}/kilometrage} & Mise \`a jour du kilom\'etrage \\
\texttt{DELETE} & \texttt{/\{id\}} & Suppression du v\'ehicule \\
\bottomrule
\end{tabularx}
\end{adrbox}

\section{Structure des donn\'ees (DTO)}

Les r\'eponses utilisent la structure suivante pour garantir la compatibilit\'e avec le schéma SQL et la passerelle :

\begin{lstlisting}[style=jsonstyle, caption={Exemple de r\'eponse VehicleResponseDto}]
{
  "id_vehicule": "550e8400-e29b-41d4-a716-446655440000",
  "immatriculation": "AB-123-CD",
  "marque": "Renault",
  "modele": "Clio",
  "annee": 2022,
  "statut": "DISPONIBLE",
  "kilometrage": 45000,
  "type": "Berline"
}
\end{lstlisting}

\section{Ports et Connectivit\'e}

Le microservice est configur\'e pour s'int\'egrer dans un environnement distribu\'e :
\begin{itemize}
    \item \textbf{Port Interne (Service)} : \texttt{4000} (configur\'e dans \texttt{application.yml})
    \item \textbf{Port Externe (Gateway/Ingress)} : \texttt{api.fleet.local} (via Traefik ou Nginx)
\end{itemize}

\section{D\'eploiement Kubernetes (K8s)}

Le service est d\'eploy\'e avec des sondes de disponibilit\'e (\textit{Liveness Probes}) configur\'ees pour assurer la r\'esilience :
\begin{itemize}
    \item \textbf{Sonde Liveness} : Appel sur \texttt{/actuator/health} (Java) ou \texttt{/health} (Python) sur le port du conteneur.
    \item \textbf{Ingress Dual-Controller} : Le fichier \texttt{ingress.yaml} d\'efinit deux ressources pour supporter \textbf{Traefik} (via \texttt{ingressClassName: traefik}) et \textbf{Nginx} (via \texttt{ingressClassName: nginx}).
\end{itemize}

\begin{infobox}
\textbf{Note sur l'acc\`es local (Windows/Minikube) :}
Pour que l'Ingress soit accessible via le nom d'h\^ote \texttt{api.fleet.local}, il est indispensable de lancer la commande \texttt{minikube tunnel} dans un terminal s\'epar\'e et de mettre \`a jour le fichier \texttt{hosts} de Windows.
\end{infobox}
% ============================================================
\chapter{Biblioth\`eques et d\'ependances}
% ============================================================

Le projet utilise \textbf{Spring Boot 3.1.5} comme framework principal.

\section{D\'ependances principales}

\begin{adrbox}{D\'ependances de production}
\begin{tabularx}{\textwidth}{l l X}
\toprule
\textbf{Biblioth\`eque} & \textbf{Artifact} & \textbf{R\^ole} \\
\midrule
Spring Boot Web & \texttt{spring-boot-starter-web} & Serveur HTTP embarqu\'e (Tomcat), API REST \\
Spring Data JPA & \texttt{spring-boot-starter-data-jpa} & ORM Hibernate, mapping entit\'e-table \\
PostgreSQL Driver & \texttt{postgresql} & Pilote JDBC pour PostgreSQL \\
Validation & \texttt{spring-boot-starter-validation} & Validation des DTOs (\texttt{@NotBlank}, \texttt{@Min}) \\
Spring Kafka & \texttt{spring-kafka} & Producteur/consommateur Kafka \\
Spring GraphQL & \texttt{spring-boot-starter-graphql} & API GraphQL avec sch\'ema SDL \\
Actuator & \texttt{spring-boot-starter-actuator} & Endpoints de sant\'e et m\'etriques \\
Micrometer & \texttt{micrometer-registry-prometheus} & Export des m\'etriques vers Prometheus \\
OpenTelemetry & \texttt{opentelemetry-spring-boot-starter} & Tracing distribu\'e \\
\bottomrule
\end{tabularx}
\end{adrbox}

\section{D\'ependances de test}

\begin{adrbox}{D\'ependances de test}
\begin{tabularx}{\textwidth}{l l X}
\toprule
\textbf{Biblioth\`eque} & \textbf{Artifact} & \textbf{R\^ole} \\
\midrule
Spring Boot Test & \texttt{spring-boot-starter-test} & JUnit 5, Mockito, MockMvc \\
H2 Database & \texttt{h2} & Base de donn\'ees en m\'emoire pour les tests \\
Spring Kafka Test & \texttt{spring-kafka-test} & Kafka embarqu\'e pour les tests \\
Spring GraphQL Test & \texttt{spring-graphql-test} & Outils de test GraphQL \\
\bottomrule
\end{tabularx}
\end{adrbox}

% ============================================================
\chapter{Mod\`ele de donn\'ees}
% ============================================================

\section{Sch\'ema SQL (PostgreSQL)}

La table \texttt{vehicule} est d\'efinie avec les contraintes suivantes :

\begin{lstlisting}[style=sqlstyle, caption={Script SQL de cr\'eation de la table vehicule}]
CREATE TYPE statut_vehicule AS ENUM (
    'DISPONIBLE', 'EN_COURSE', 'EN_MAINTENANCE'
);

CREATE TABLE vehicule (
    id_vehicule     UUID PRIMARY KEY DEFAULT gen_random_uuid(),
    immatriculation VARCHAR(20)      NOT NULL UNIQUE,
    marque          VARCHAR(100)     NOT NULL,
    modele          VARCHAR(100)     NOT NULL,
    annee           SMALLINT         NOT NULL
                    CHECK (annee >= 1900
                    AND annee <= EXTRACT(YEAR FROM NOW()) + 1),
    statut          statut_vehicule  NOT NULL DEFAULT 'DISPONIBLE',
    kilometrage     INTEGER          NOT NULL DEFAULT 0
                    CHECK (kilometrage >= 0),
    type            VARCHAR(50)      NOT NULL,
    created_at      TIMESTAMPTZ      NOT NULL DEFAULT NOW(),
    updated_at      TIMESTAMPTZ
);
\end{lstlisting}

\begin{infobox}
\textbf{Points cl\'es du sch\'ema :}
\begin{itemize}[nosep]
    \item \textbf{UUID} comme cl\'e primaire, g\'en\'er\'e via \texttt{gen\_random\_uuid()}
    \item \textbf{Enum PostgreSQL} pour le statut, garantissant l'int\'egrit\'e
    \item \textbf{Contraintes CHECK} sur l'ann\'ee et le kilom\'etrage
    \item \textbf{Unicit\'e} de l'immatriculation
    \item \textbf{Index} sur statut et immatriculation
\end{itemize}
\end{infobox}

\section{Entit\'e JPA}

L'entit\'e Java \texttt{Vehicle} est mapp\'ee sur la table SQL :

\begin{lstlisting}[style=javastyle, language=Java, caption={Extrait de Vehicle.java}]
@Entity
@Table(name = "vehicule")
public class Vehicle {

    @Id
    @GeneratedValue(strategy = GenerationType.UUID)
    @Column(name = "id_vehicule")
    private UUID id;

    @Column(name = "immatriculation", nullable = false,
            unique = true, length = 20)
    private String immatriculation;

    @Column(nullable = false, length = 100)
    private String marque;

    @Column(nullable = false, length = 100)
    private String modele;

    @Column(name = "annee", nullable = false)
    private Short annee;

    @Column(nullable = false)
    @Enumerated(EnumType.STRING)
    private VehicleStatus statut = VehicleStatus.DISPONIBLE;

    @Column(nullable = false)
    private Integer kilometrage = 0;

    public enum VehicleStatus {
        DISPONIBLE, EN_COURSE, EN_MAINTENANCE
    }

    @PrePersist
    protected void onCreate() {
        this.createdAt = LocalDateTime.now();
        this.updatedAt = LocalDateTime.now();
    }

    @PreUpdate
    protected void onUpdate() {
        this.updatedAt = LocalDateTime.now();
    }
}
\end{lstlisting}

\section{Correspondance entit\'e / table}

\begin{adrbox}{Mapping JPA --- PostgreSQL}
\begin{tabular}{l l l}
\toprule
\textbf{Champ Java} & \textbf{Colonne SQL} & \textbf{Type SQL} \\
\midrule
\texttt{id} (UUID) & \texttt{id\_vehicule} & \texttt{UUID} \\
\texttt{immatriculation} & \texttt{immatriculation} & \texttt{VARCHAR(20)} \\
\texttt{marque} & \texttt{marque} & \texttt{VARCHAR(100)} \\
\texttt{modele} & \texttt{modele} & \texttt{VARCHAR(100)} \\
\texttt{annee} (Short) & \texttt{annee} & \texttt{SMALLINT} \\
\texttt{statut} (Enum) & \texttt{statut} & \texttt{statut\_vehicule} \\
\texttt{kilometrage} & \texttt{kilometrage} & \texttt{INTEGER} \\
\texttt{type} & \texttt{type} & \texttt{VARCHAR(50)} \\
\texttt{createdAt} & \texttt{created\_at} & \texttt{TIMESTAMPTZ} \\
\texttt{updatedAt} & \texttt{updated\_at} & \texttt{TIMESTAMPTZ} \\
\bottomrule
\end{tabular}
\end{adrbox}

% ============================================================
\chapter{API REST --- CRUD}
% ============================================================

Le contr\^oleur REST expose les op\'erations CRUD classiques via \texttt{VehicleController}.

\section{Endpoints disponibles}

\begin{adrbox}{Endpoints REST du service v\'ehicule}
\begin{tabularx}{\textwidth}{l l l X}
\toprule
\textbf{M\'ethode} & \textbf{URL} & \textbf{Action} & \textbf{Description} \\
\midrule
POST & \texttt{/api/vehicles} & Create & Cr\'eer un nouveau v\'ehicule \\
GET & \texttt{/api/vehicles} & Read All & Lister tous les v\'ehicules \\
GET & \texttt{/api/vehicles/\{id\}} & Read One & Obtenir un v\'ehicule par UUID \\
GET & \texttt{/api/vehicles/statut/\{s\}} & Filter & Filtrer par statut \\
PUT & \texttt{/api/vehicles/\{id\}} & Update & Mettre \`a jour un v\'ehicule \\
DELETE & \texttt{/api/vehicles/\{id\}} & Delete & Supprimer un v\'ehicule \\
\bottomrule
\end{tabularx}
\end{adrbox}

\section{Contr\^oleur REST}

\begin{lstlisting}[style=javastyle, language=Java, caption={VehicleController.java}]
@RestController
@RequestMapping("/api/vehicles")
public class VehicleController {

    private final VehicleService vehicleService;

    @PostMapping
    public ResponseEntity<VehicleResponseDto> createVehicle(
            @Valid @RequestBody VehicleRequestDto request) {
        VehicleResponseDto response =
            vehicleService.createVehicle(request);
        return ResponseEntity.status(HttpStatus.CREATED)
                             .body(response);
    }

    @GetMapping("/{id}")
    public ResponseEntity<VehicleResponseDto> getVehicleById(
            @PathVariable UUID id) {
        VehicleResponseDto vehicle =
            vehicleService.getVehicleById(id);
        return ResponseEntity.ok(vehicle);
    }

    @PutMapping("/{id}")
    public ResponseEntity<VehicleResponseDto> updateVehicle(
            @PathVariable UUID id,
            @Valid @RequestBody VehicleRequestDto request) {
        VehicleResponseDto response =
            vehicleService.updateVehicle(id, request);
        return ResponseEntity.ok(response);
    }

    @DeleteMapping("/{id}")
    public ResponseEntity<Void> deleteVehicle(
            @PathVariable UUID id) {
        vehicleService.deleteVehicle(id);
        return ResponseEntity.noContent().build();
    }
}
\end{lstlisting}

\section{DTOs (Data Transfer Objects)}

Le projet utilise deux DTOs pour d\'ecoupler l'API de l'entit\'e interne :

\begin{itemize}
    \item \textbf{VehicleRequestDto} --- Donn\'ees entrantes avec validation :
    \begin{itemize}[nosep]
        \item \texttt{@NotBlank} sur marque, mod\`ele, immatriculation, type
        \item \texttt{@NotNull} sur l'ann\'ee
        \item \texttt{@Min(0)} sur le kilom\'etrage
    \end{itemize}
    \item \textbf{VehicleResponseDto} --- Donn\'ees sortantes avec ID UUID, timestamps, etc.
\end{itemize}

\section{Couche Service}

\begin{lstlisting}[style=javastyle, language=Java, caption={Extrait de VehicleServiceImpl.java}]
@Service
@Transactional
public class VehicleServiceImpl implements VehicleService {

    private final VehicleRepository vehicleRepository;
    private final VehicleProducer vehicleProducer;

    @Override
    public VehicleResponseDto createVehicle(
            VehicleRequestDto request) {
        // Verification d'unicite de l'immatriculation
        if (vehicleRepository.existsByImmatriculation(
                request.getImmatriculation())) {
            throw new IllegalArgumentException(
                "Immatriculation deja existante");
        }
        Vehicle vehicle = mapToEntity(request);
        Vehicle saved = vehicleRepository.save(vehicle);

        // Publication de l'evenement Kafka
        vehicleProducer.sendVehicleEvent(
            "VEHICLE_CREATED", mapToDto(saved));
        return mapToDto(saved);
    }
}
\end{lstlisting}

\begin{infobox}
\textbf{Points importants de la couche Service :}
\begin{itemize}[nosep]
    \item \texttt{@Transactional} assure l'int\'egrit\'e des op\'erations
    \item \texttt{@Transactional(readOnly = true)} optimise les lectures
    \item V\'erifications m\'etier (unicit\'e d'immatriculation)
    \item Publication automatique d'\'ev\'enements Kafka apr\`es chaque op\'eration
\end{itemize}
\end{infobox}

% ============================================================
\chapter{Apache Kafka --- Communication asynchrone}
% ============================================================

\section{R\^ole de Kafka dans l'architecture}

Apache Kafka est utilis\'e comme syst\`eme de messagerie asynchrone pour :
\begin{itemize}
    \item \textbf{D\'ecoupler} les services --- le producteur n'attend pas de r\'eponse
    \item \textbf{Publier des \'ev\'enements} sur le cycle de vie des v\'ehicules
    \item \textbf{Permettre le replay} des \'ev\'enements en cas de besoin
    \item \textbf{Alimenter d'autres services} (notifications, analytics, etc.)
\end{itemize}

\section{Producteur Kafka}

\begin{lstlisting}[style=javastyle, language=Java, caption={VehicleProducer.java}]
@Service
public class VehicleProducer {

    private final KafkaTemplate<String, String> kafkaTemplate;
    private final ObjectMapper objectMapper;

    @Value("${app.kafka.topic.vehicle-events:vehicle-events}")
    private String vehicleEventsTopic;

    public void sendVehicleEvent(String eventType,
                                  VehicleResponseDto vehicle) {
        String key = eventType + "-" + vehicle.getId();
        String payload = objectMapper.writeValueAsString(
            new VehicleEvent(eventType, vehicle));
        kafkaTemplate.send(vehicleEventsTopic, key, payload);
    }

    // Classe interne -- structure de l'evenement
    public static class VehicleEvent {
        private String eventType;
        private VehicleResponseDto vehicle;
        private long timestamp;
    }
}
\end{lstlisting}

\begin{adrbox}{Types d'\'ev\'enements Kafka}
\begin{tabular}{l l}
\toprule
\textbf{\'Ev\'enement} & \textbf{D\'eclencheur} \\
\midrule
\texttt{VEHICLE\_CREATED} & Cr\'eation d'un v\'ehicule \\
\texttt{VEHICLE\_UPDATED} & Mise \`a jour d'un v\'ehicule \\
\texttt{VEHICLE\_DELETED} & Suppression d'un v\'ehicule \\
\bottomrule
\end{tabular}
\end{adrbox}

\section{Consommateur Kafka}

\begin{lstlisting}[style=javastyle, language=Java, caption={VehicleConsumer.java}]
@Service
public class VehicleConsumer {

    @KafkaListener(
        topics = "${app.kafka.topic.vehicle-events}",
        groupId = "${spring.kafka.consumer.group-id}")
    public void consumeVehicleEvent(String message) {
        log.info("Evenement vehicule recu: {}", message);
        // Traitement : notifications, mises a jour, etc.
    }
}
\end{lstlisting}

\section{Configuration Kafka}

\begin{lstlisting}[style=yamlstyle, caption={Configuration Kafka dans application.yml}]
spring:
  kafka:
    bootstrap-servers: localhost:9092
    producer:
      key-serializer: org.apache.kafka.common.serialization
                      .StringSerializer
      value-serializer: org.apache.kafka.common.serialization
                        .StringSerializer
    consumer:
      group-id: vehicule-service-group
      auto-offset-reset: earliest
\end{lstlisting}

% ============================================================
\chapter{GraphQL --- API flexible}
% ============================================================

\section{Pourquoi GraphQL ?}

GraphQL offre des avantages compl\'ementaires \`a REST :
\begin{itemize}
    \item \textbf{Requ\^etes flexibles} --- Le client choisit exactement les champs voulus
    \item \textbf{Un seul endpoint} --- \texttt{/graphql} remplace plusieurs endpoints REST
    \item \textbf{Pas de sur-r\'ecup\'eration} --- Contrairement \`a REST qui retourne tous les champs
    \item \textbf{Introspection} --- Documentation automatique du sch\'ema
\end{itemize}

\section{Sch\'ema GraphQL}

\begin{lstlisting}[style=graphqlstyle, caption={schema.graphqls}]
type Vehicle {
    id: ID!
    marque: String!
    modele: String!
    immatriculation: String!
    statut: String!
    kilometrage: Int
    type: String!
    annee: Int!
    createdAt: String
    updatedAt: String
}

type Query {
    allVehicles: [Vehicle!]!
    vehicleById(id: ID!): Vehicle
    vehiclesByStatut(statut: String!): [Vehicle!]!
}

type Mutation {
    createVehicle(
        marque: String!, modele: String!,
        immatriculation: String!, type: String!,
        statut: String, kilometrage: Int, annee: Int!
    ): Vehicle!
    updateVehicle(
        id: ID!, marque: String!, modele: String!,
        immatriculation: String!, type: String!,
        statut: String, kilometrage: Int, annee: Int!
    ): Vehicle!
    deleteVehicle(id: ID!): Boolean!
}
\end{lstlisting}

\section{Contr\^oleur GraphQL}

\begin{lstlisting}[style=javastyle, language=Java, caption={Extrait de VehicleGraphQLController.java}]
@Controller
public class VehicleGraphQLController {

    private final VehicleService vehicleService;

    @QueryMapping
    public List<VehicleResponseDto> allVehicles() {
        return vehicleService.getAllVehicles();
    }

    @QueryMapping
    public VehicleResponseDto vehicleById(@Argument String id) {
        return vehicleService.getVehicleById(
            UUID.fromString(id));
    }

    @MutationMapping
    public VehicleResponseDto createVehicle(
            @Argument String marque,
            @Argument String modele,
            @Argument String immatriculation,
            @Argument String type,
            @Argument Integer annee) {
        VehicleRequestDto request = new VehicleRequestDto();
        request.setMarque(marque);
        request.setModele(modele);
        request.setImmatriculation(immatriculation);
        request.setType(type);
        request.setAnnee(annee != null
            ? annee.shortValue() : null);
        return vehicleService.createVehicle(request);
    }
}
\end{lstlisting}

\section{Exemple de requ\^ete GraphQL}

\begin{lstlisting}[style=graphqlstyle, caption={Requ\^ete pour obtenir marque et mod\`ele}]
query {
  allVehicles {
    marque
    modele
    statut
  }
}
\end{lstlisting}

\begin{infobox}
L'interface \textbf{GraphiQL} est activ\'ee et accessible \`a \texttt{http://localhost:4000/graphiql} pour tester les requ\^etes de mani\`ere interactive.
\end{infobox}

% ============================================================
\chapter{Tests automatis\'es}
% ============================================================

\section{Strat\'egie de test}

Le projet impl\'emente deux niveaux de tests :

\begin{enumerate}
    \item \textbf{Tests unitaires} --- Couche service avec Mockito
    \item \textbf{Tests d'int\'egration} --- Couche contr\^oleur avec MockMvc
\end{enumerate}

\section{Tests unitaires de la couche Service}

Les tests de \texttt{VehicleServiceTest} utilisent \textbf{Mockito} pour isoler la logique m\'etier :

\begin{lstlisting}[style=javastyle, language=Java, caption={Extrait de VehicleServiceTest.java}]
@ExtendWith(MockitoExtension.class)
class VehicleServiceTest {

    @Mock
    private VehicleRepository vehicleRepository;
    @Mock
    private VehicleProducer vehicleProducer;
    @InjectMocks
    private VehicleServiceImpl vehicleService;

    @Test
    void createVehicle_Success() {
        when(vehicleRepository.existsByImmatriculation(
            "AB-123-CD")).thenReturn(false);
        when(vehicleRepository.save(any(Vehicle.class)))
            .thenReturn(vehicle);

        VehicleResponseDto result =
            vehicleService.createVehicle(requestDto);

        assertNotNull(result);
        assertEquals("Renault", result.getMarque());
        verify(vehicleProducer).sendVehicleEvent(
            eq("VEHICLE_CREATED"),
            any(VehicleResponseDto.class));
    }

    @Test
    void createVehicle_DuplicateImmatriculation() {
        when(vehicleRepository.existsByImmatriculation(
            "AB-123-CD")).thenReturn(true);

        assertThrows(IllegalArgumentException.class,
            () -> vehicleService.createVehicle(requestDto));
        verify(vehicleRepository, never())
            .save(any(Vehicle.class));
    }
}
\end{lstlisting}

\section{Tests d'int\'egration du contr\^oleur REST}

\begin{lstlisting}[style=javastyle, language=Java, caption={Extrait de VehicleControllerTest.java}]
@WebMvcTest(VehicleController.class)
class VehicleControllerTest {

    @Autowired
    private MockMvc mockMvc;
    @MockBean
    private VehicleService vehicleService;

    @Test
    void createVehicle_ReturnsCreated() throws Exception {
        when(vehicleService.createVehicle(
            any(VehicleRequestDto.class)))
            .thenReturn(responseDto);

        mockMvc.perform(post("/api/vehicles")
                .contentType(MediaType.APPLICATION_JSON)
                .content(objectMapper
                    .writeValueAsString(requestDto)))
            .andExpect(status().isCreated())
            .andExpect(jsonPath("$.marque").value("Renault"))
            .andExpect(jsonPath("$.immatriculation")
                .value("AB-123-CD"));
    }

    @Test
    void createVehicle_ValidationError() throws Exception {
        VehicleRequestDto invalid = new VehicleRequestDto();

        mockMvc.perform(post("/api/vehicles")
                .contentType(MediaType.APPLICATION_JSON)
                .content(objectMapper
                    .writeValueAsString(invalid)))
            .andExpect(status().isBadRequest());
    }
}
\end{lstlisting}

\section{Cas de test couverts}

\begin{adrbox}{Couverture des tests}
\begin{tabularx}{\textwidth}{l l X}
\toprule
\textbf{Classe} & \textbf{Test} & \textbf{Sc\'enario} \\
\midrule
Service & createVehicle\_Success & Cr\'eation r\'eussie + \'ev\'enement Kafka \\
Service & createVehicle\_Duplicate & Rejet si immatriculation dupliqu\'ee \\
Service & getVehicleById\_Success & R\'ecup\'eration par UUID \\
Service & getVehicleById\_NotFound & Exception si inexistant \\
Service & getAllVehicles\_Success & Liste compl\`ete \\
Service & getVehiclesByStatut & Filtrage par statut \\
Service & updateVehicle\_Success & Mise \`a jour + \'ev\'enement Kafka \\
Service & updateVehicle\_NotFound & Exception si inexistant \\
Service & deleteVehicle\_Success & Suppression + \'ev\'enement Kafka \\
Service & deleteVehicle\_NotFound & Exception si inexistant \\
\midrule
Controller & createVehicle\_Created & HTTP 201 + JSON \\
Controller & getAllVehicles\_List & HTTP 200 + tableau \\
Controller & getVehicleById & HTTP 200 + d\'etail \\
Controller & updateVehicle\_Updated & HTTP 200 + modifi\'e \\
Controller & deleteVehicle\_NoContent & HTTP 204 \\
Controller & validation\_BadRequest & HTTP 400 si invalide \\
\bottomrule
\end{tabularx}
\end{adrbox}

% ============================================================
\chapter{Observabilit\'e --- OpenTelemetry SDK}
% ============================================================

L'impl\'ementation de l'observabilit\'e repose sur les \textbf{3 piliers d'OpenTelemetry} (Traces, M\'etriques, Logs), int\'egr\'es via le SDK Java.

\section{Le SDK et la configuration centralis\'ee}

Une classe \texttt{VehicleTelemetry} centralise l'acc\`es au \textbf{Tracer} (pour les traces) et au \textbf{Meter} (pour les m\'etriques).

\begin{lstlisting}[style=javastyle, caption={Extraits de VehicleTelemetry.java}]
@Component
public class VehicleTelemetry {
    private final Tracer tracer;
    private final LongCounter vehiclesCreatedCounter;
    private final LongHistogram operationDurationHistogram;

    public VehicleTelemetry(OpenTelemetry openTelemetry) {
        this.tracer = openTelemetry.getTracer("vehicule-service");
        Meter meter = openTelemetry.getMeter("vehicule-service");
        
        // Compteur : nombre total de creations
        this.vehiclesCreatedCounter = meter.counterBuilder("vehicles.created.total").build();
        
        // Histogramme : duree des operations
        this.operationDurationHistogram = meter.histogramBuilder("vehicles.operation.duration").ofLongs().build();
    }
}
\end{lstlisting}

\section{Pilier 1 : Traces Distribu\'ees}

Les traces permettent de suivre une requ\^ete. Nous utilisons deux m\'ethodes :
\begin{itemize}
    \item \textbf{Instrumentation d\'eclarative} : Utilisation de \texttt{@WithSpan} pour cr\'eer automatiquement un span sur une m\'ethode.
    \item \textbf{Attributs personnalis\'es} : Utilisation de \texttt{@SpanAttribute} pour injecter des param\`etres (ex: UUID du v\'ehicule) dans la trace.
\end{itemize}

\begin{lstlisting}[style=javastyle, caption={Traces dans la couche Service}]
@WithSpan("createVehicle")
public VehicleResponseDto createVehicle(VehicleRequestDto request) {
    Span.current().setAttribute("vehicle.marque", request.getMarque());
    // ... logique metier ...
}
\end{lstlisting}

\section{Pilier 2 : M\'etriques Personnalis\'ees}

Contrairement aux traces, les m\'etriques sont des donn\'ees num\'eriques agr\'eg\'ees. Nous mesurons :
\begin{itemize}
    \item Le nombre total de v\'ehicules cr\'e\'es et supprim\'es.
    \item Le taux d'erreurs par type d'op\'eration.
    \item La distribution des temps de traitement (latence).
\end{itemize}

\section{Pilier 3 : Logs Structur\'es et Corr\'elation}

Pour relier les logs aux traces, nous injectons le \texttt{traceId} et le \texttt{spanId} dans le contexte de diagnostic (\textbf{MDC}).

\begin{lstlisting}[style=xmlstyle, caption={Configuration Logback (logback-spring.xml)}]
<encoder>
    <pattern>
        %d %-5level [traceId=%X{traceId:-none}, spanId=%X{spanId:-none}] %logger - %msg%n
    </pattern>
</encoder>
\end{lstlisting}

\section{V\'erification du fonctionnement}

Pour s'assurer que l'observabilit\'e est active :

\begin{adrbox}{Plan de v\'erification}
\begin{enumerate}
    \item \textbf{Logs} : Lancer l'application et effectuer un appel API. Le log doit afficher \texttt{traceId=...} avec un ID hexad\'ecimal au lieu de \texttt{none}.
    \item \textbf{M\'etriques} : Acc\'eder \`a \texttt{/actuator/prometheus}. Rechercher la ligne \texttt{vehicles\_created\_total} pour voir le compteur augmenter.
    \item \textbf{Traces} : Si un collecteur (Jaeger/Tempo) est configur\'e, les spans \texttt{createVehicle} appara\^\i tront avec leurs attributs (\texttt{vehicle.marque}).
\end{enumerate}
\end{adrbox}


% ============================================================
\chapter{Configuration de l'application}
% ============================================================

\begin{lstlisting}[style=yamlstyle, caption={application.yml -- Configuration compl\`ete}]
server:
  port: 4000

spring:
  application:
    name: vehicule-service
  datasource:
    url: jdbc:postgresql://localhost:5432/fleet_vehicles
    username: postgres
    password: postgres
  jpa:
    hibernate:
      ddl-auto: update
    properties:
      hibernate:
        dialect: org.hibernate.dialect.PostgreSQLDialect
  graphql:
    graphiql:
      enabled: true
    path: /graphql

management:
  endpoints:
    web:
      exposure:
        include: health,info,metrics,prometheus
\end{lstlisting}

\begin{infobox}
\textbf{Note :} Le mode \texttt{ddl-auto: update} permet \`a Hibernate de mettre \`a jour le sch\'ema automatiquement au d\'emarrage. En production, il est recommand\'e d'utiliser un outil de migration comme \textbf{Flyway} ou \textbf{Liquibase}.
\end{infobox}

% ============================================================
\chapter{Conclusion}
% ============================================================

Le microservice \textbf{vehicule-service} impl\'emente une architecture moderne et compl\`ete :

\begin{itemize}
    \item \textbf{CRUD REST} structur\'e avec validation, DTOs et gestion d'erreurs
    \item \textbf{GraphQL} offrant une API flexible compl\'ementaire
    \item \textbf{Kafka} pour la communication \'ev\'enementielle asynchrone inter-services
    \item \textbf{Tests automatis\'es} couvrant la logique m\'etier et les endpoints HTTP
    \item \textbf{Observabilit\'e} via OpenTelemetry, Actuator et Prometheus
    \item \textbf{Coh\'erence entit\'e/base} avec un mapping JPA align\'e sur le sch\'ema PostgreSQL
\end{itemize}

Cette architecture permet au service d'\^etre \textit{scalable}, \textit{maintenable} et \textit{observable} dans un contexte de d\'eploiement microservices.

\end{document}
